\scanpage{173}

groupes par les données suivantes

\[
\left\{\begin{aligned}
\text{a)}\quad&\text{Un diagramme de groupes}\\
 &\widetilde U\subset\widetilde B''\subset G_0,\quad\text{avec $\widetilde U$ distingué dans $\widetilde B''$}\\
 &\deleted{\unsure{b)}}\quad\text{d'où un groupe quotient}\\
 &Z=\widetilde B''/\widetilde U\qquad(\widetilde B''\text{ ext. de $Z$ par $\widetilde U$})\\
\text{b)}\quad&\text{Un \deleted{hom.} hom de groupes }Z\xrightarrow{\unsure{\epsilon|\varepsilon}}Z'\quad\text{et}\\
 &\text{une extension \inserted{$B'$} de $Z'$ par le $\widetilde U$ du a,}\\
 &\text{de telle façon que $\widetilde B''$ soit déduit de}\\
 &\text{cette extension par l'hom $Z\mathrel{\text{---}}Z'$, i.e.}\\
 &\text{qu'on ait un hom de suites exactes (46)}\\
\text{c)}\quad&\text{Une opération de $B'$ sur $G_0$, de telle}\\
 &\text{façon que l'opération composée}\\
 &\widetilde B''\xrightarrow{\psi}B'\mathrel{\text{---}}\operatorname{Aut}_{\rm gr}(G_0)\\
 &\text{soit l'opération tautologique déduite du}\\
 &\text{plongement a) de $\widetilde B''$ dans $G_0$.}
\end{aligned}\right.\tag*{(50)}
\]

% slanted
\margin{oublié : la condition que pour $g\in\widetilde B^{\unsure{\prime\prime|\prime}}$, $\Theta(g)|U=\operatorname{int}(g)|U$, et que $\unsure{\Theta g)|\Theta(g)|\theta(g)}$ normalise $U$ et centralise $\widetilde B''$ mod $U$}

En termes des données (26), on a

\[
\left\{\begin{aligned}
 &\deleted{\unsure{B}}\widetilde U=G_{0s_0}\underset{\unsure{\mathrm{df}|\times}}{=}\deleted{\unsure{B}}\{g\in G_0\mid g.s_0=s_0\}\\
 &\widetilde B''=\{g\in G_0\mid gs_0\in Z.s\}\\
 &Z'=G/G_0,\quad Z\longrightarrow Z'\quad\text{le composé }Z\underset{\rm incl.}{\subset}G\xrightarrow{\rm can}\underbrace{G/G_0}_{Z'}\\
 &B'=\deleted{\unsure{\{g\in}}G_{s_0}=\{g\in G\mid gs_0=s_0\}\quad\text{de sorte que}\\
 &\deleted{\unsure{\widetilde U=B'|Z=}}\quad\widetilde U=B'\cap G_0,\quad B'\longrightarrow Z'\ \text{composé }B'\underset{\rm incl}{\subset}G\xrightarrow{\rm can}\underbrace{G/G_0}_{Z'}\\
 &\text{Op. de $B'$ sur $G_0$ induite par celle de $G$ sur}\\
 &\text{le sous-groupe distingué $G_0$}
\end{aligned}\right.\tag*{(51)}
\]

En termes des données (\unsure{50|80}), on récupère (26) par

\[
\left\{\begin{aligned}
\text{a)}\quad&G\ \text{ext. de $Z'$ par $G_0$ déduite de \deleted{$\unsure{B'}$} l'ext \inserted{b)} $B'$ de $Z'$ par $\widetilde U$,}\\
 &\text{le plongement $\widetilde U\mathrel{\lhook\joinrel-}G_0$, et l'opération \inserted{c)} de $B'$ sur $G_0$.}\\
\text{b)}\quad&S\simeq G_0/\widetilde U\simeq G/B'\\
\text{c)}\quad&Z=\operatorname{Im}(\widetilde B''\xrightarrow{\unsure{\varphi|\phi}}G),\quad\varphi(g)=g.\psi(g)^{-1}\\
 &\hspace{16em}=\psi(g)^{-1}.g
\end{aligned}\right.\tag*{(52)}
\]

\margin{produit dans $G$ qui contient $\widetilde B'',B'$, $\psi:\widetilde B''\mathrel{\text{---}}B'$ défini dans (50 c).}

\margin{\uline{NB} L'hom $\varphi:\widetilde B''\longrightarrow G$ se factorise par \unsure{le|son} quotient $\widetilde B''/\widetilde U\simeq Z$}

% archive sheet number: 172
